\section{Uniform smoothing of matching averages}\label{sec:matching-smoothing}

This appendix proves an independent smoothing theorem for averages over
perfect matchings with prescribed numbers of edges between vertex cells.
The number of cells is fixed, but their sizes and the prescribed counts
may vary arbitrarily.  The conclusion holds even when the prescribed
table has very small probability.  Besides an operator-norm estimate,
we obtain a single polynomial that describes all the low-degree averages
at a given partition.

The proof uses the same comparison between a fourth moment and eigenspace
dimensions as Lemma~\ref{lem:local-averaging}.  Here the inputs are
functions of perfect matchings, and removing products of a bounded number
of edge indicators makes the relevant dimensions grow by an arbitrarily
large power of the number of vertices.

Fix a positive integer \(\ell\).  For even \(N\), let \(\mathcal X\)
be the perfect matchings on a fixed \(N\)-element vertex set, and let
\(\mathcal Z\) be the ordered partitions into cells of positive sizes
\(s_1,\ldots,s_\ell\), with \(\sum_us_u=N\).  Both spaces carry the
uniform probability measure.  The \emph{table} of a matching at a
partition records its internal edge counts \(t_{uu}\) and its edge
counts \(t_{uv}=t_{vu}\) between distinct cells.  A table is admissible
if its entries are nonnegative integers and
\begin{equation}\label{smooth:eq-table-degrees}
  2t_{uu}+\sum_{v\ne u}t_{uv}=s_u
  \qquad (1\le u\le\ell).
\end{equation}
These conditions suffice: allocate the required endpoints in each
cell, pair the allocated endpoints between cells, and match the
remaining vertices internally.

For an admissible table, define
\(T_t:L^2(\mathcal X)\longrightarrow L^2(\mathcal Z)\) by
\[
  (T_tf)(Z)=\E[\,f(M)\mid M\text{ has table }t\text{ at }Z\,].
\]
Let \(\mathcal P_d\) be the span of products of at most \(d\) edge
indicators \(x_e(M)=\mathbf1_{\{e\in M\}}\), including the constant,
and let \(\Pi_d\) be its orthogonal projection.

\begin{theorem}[Uniform smoothing]\label{thm:matching-smoothing}
For every fixed \(\ell\ge1\) and \(K>0\), there are an integer
\(d\ge1\) and \(N_0\) such that, for every even \(N\ge N_0\), every
positive integer cell-size vector summing to \(N\), and every admissible table,
\begin{equation}\label{smooth:eq-smoothing}
  \normop{T_t(I-\Pi_d)}\le N^{-K}.
\end{equation}
Moreover, for each fixed partition \(Z\) and function \(f\), the values
\((T_t\Pi_df)(Z)\), as \(t\) varies over all admissible tables, are
given by a single polynomial of total degree at most \(d\) in the
intercell counts \((t_{uv})_{u<v}\).
\end{theorem}

The approximation error in \eqref{smooth:eq-smoothing} is measured in
\(L^2(\mathcal Z)\), with respect to the uniform choice of partition.
We prove the analytic estimate by combining a polynomial bound on a
fourth moment with a lower bound on the dimensions of the permutation
representations outside \(\mathcal P_d\).  The polynomial assertion will
then follow from a direct containment count.

\subsection{The fourth moment}

\begin{lemma}\label{smooth:lem-fourth-moment}
Put \(a_\ell=\ell(\ell+1)+1\) and
\(b_\ell=2a_\ell+\ell^2\).  For every admissible table,
\begin{equation}\label{smooth:eq-fourth-moment}
  \tr\bigl((T_t^*T_t)^2\bigr)
  \le e^{2a_\ell}(N+1)^{b_\ell}.
\end{equation}
\end{lemma}

\begin{proof}
Write \(h=N/2\).  For a fixed matching, the number of partitions
giving it table \(t\) is
\[
 D_t=\frac{h!}{\prod_{u\le v}t_{uv}!}\,
                  2^{\sum_{u<v}t_{uv}},
 \qquad
 |\mathcal Z|=\frac{N!}{\prod_us_u!}.
\]
The first formula assigns cell types to the \(h\) edges and chooses
which endpoint occupies which cell for each intercell edge.
Permutation transitivity shows that the table probability is
\(p_t=D_t/|\mathcal Z|\), whether the matching or the partition is
fixed.  Thus the density kernel of \(T_t\), relative to the uniform
measures, is the table indicator divided by \(p_t\).

Let \(J(M,M')\) count partitions giving table \(t\) to both matchings.
The density kernel of \(\mathcal A=T_t^*T_t\) is consequently
\begin{equation}\label{smooth:eq-matching-kernel}
 \kappa(M,M')=\frac{|\mathcal Z|}{D_t^2}J(M,M'),
 \qquad
 \tr(\mathcal A^2)=\E_{M,M'}\kappa(M,M')^2,
\end{equation}
where \(M,M'\) are independent and uniform.

We first bound this kernel using the alternating cycles of
\(M\cup M'\).  Define
\[
 \pi_u=\frac{s_u}{N},\qquad
 \gamma_{uu}=\frac{2t_{uu}}N,\qquad
 \gamma_{uv}=\frac{t_{uv}}N\quad(u\ne v).
\]
The array \(\gamma\) is symmetric, has total sum one, and has row sums
\(\pi\).  For a probability vector or array write
\(H(p)=-\sum_jp_j\log p_j\), with \(0\log0=0\).  The factorial
formulas imply
\begin{equation}\label{smooth:eq-table-entropy}
 \log\frac{D_t^2}{|\mathcal Z|}
 =N\bigl(H(\gamma)-H(\pi)\bigr)+\Delta,
 \qquad
 \Delta\ge-a_\ell\bigl(1+\log(N+1)\bigr).
\end{equation}
To check the remainder uniformly, including zero table entries, put
\(R_0=0\) and
\(R_j=\log(j!)-j\log j+j\) for \(j\ge1\).
Integral comparison gives \(0\le R_j\le1+\log(j+1)\), and direct
substitution gives
\[
 \Delta=2R_h-2\sum_{u\le v}R_{t_{uv}}-R_N+\sum_uR_{s_u}.
\]
There are \(\ell(\ell+1)/2\) unordered table entries, proving the
stated lower bound.  In the leading term, splitting each unordered
off-diagonal frequency into two ordered frequencies accounts for
the power of two in \(D_t\).

Retain doubled edges in \(M\cup M'\), so that its components are
alternating cycles, including cycles of length two.  Orient each
cycle and write \(c=c(M,M')\) for their number.  The matrix
\(W_{uv}=\gamma_{uv}/\pi_u\) is nonnegative and stochastic.  Give a
vertex coloring by \(1,\ldots,\ell\) the weight obtained by
multiplying \(W_{uv}\) over the oriented edges, where \(u,v\) are
the tail and head colors.

Every coloring counted by \(J(M,M')\) has weight
\begin{equation}\label{smooth:eq-table-color-weight}
 \exp\{-N(H(\gamma)-H(\pi))\}.
\end{equation}
Indeed, each vertex occurs once as a tail, so the denominator is
\(\prod_u\pi_u^{s_u}\).  The two matchings together contribute
\(2t_{uv}\) edges of each unordered type.  Symmetry of \(\gamma\)
makes their orientations irrelevant, and the numerator is
\(\prod_{u\le v}\gamma_{uv}^{2t_{uv}}\), with zero exponents omitted.
These expressions give \eqref{smooth:eq-table-color-weight}; no counted
coloring uses a zero entry of \(\gamma\).

For a cycle of length \(a\), the sum of weights over all its colorings
is \(\tr(W^a)\le\ell\).  Summing over all vertex colorings therefore
bounds the total weight of the colorings counted by \(J\) by
\(\ell^c\).  Equations \eqref{smooth:eq-matching-kernel}--\eqref{smooth:eq-table-color-weight}
yield
\begin{equation}\label{smooth:eq-matching-kernel-bound}
 \kappa(M,M')\le e^{a_\ell}(N+1)^{a_\ell}\ell^{c(M,M')}.
\end{equation}

It remains to average the cycle factor.  Set \(b=\ell^2\).  The
number \(b^c\) counts colorings that are constant on every edge of
both matchings.  Fix \(M\), and suppose its \(h\) edges have color
multiplicities \(j_1,\ldots,j_b\), summing to \(h\).  Averaging the
number of such colorings over \(M'\) gives the contribution
\[
 \frac{h!}{\prod_i j_i!}
 \frac{\prod_i(2j_i-1)!!}{(2h-1)!!}
 =\frac{\prod_i\binom{2j_i}{j_i}}{\binom{2h}{h}}
 \le1,
\]
where \((-1)!!=1\).  For the inequality, partition a \(2h\)-element
set into blocks of sizes \(2j_i\): the numerator counts only those
\(h\)-element subsets taking \(j_i\) elements from each block.
There are at most \((h+1)^b\) possible multiplicity vectors.  Hence
\(\E\ell^{2c}\le(h+1)^{\ell^2}\).  Squaring
\eqref{smooth:eq-matching-kernel-bound} and using
\eqref{smooth:eq-matching-kernel} proves \eqref{smooth:eq-fourth-moment}.
\end{proof}

\subsection{Dimensions outside the edge-product span}

The fourth-moment bound alone does not make the operator small.
The additional fact is that each eigenspace of the compression of
\(T_t^*T_t\) to \(\mathcal P_d^\perp\) is invariant, and every nonzero
such space has large dimension.  Here we use the standard
representation theory of the symmetric group over the complex
numbers: irreducibles are indexed by Young diagrams, restriction
from \(S_N\) to \(S_{N-1}\) deletes a corner box, and dimensions count
standard Young tableaux.  The one-row and one-column diagrams give
the trivial and sign representations.  These facts, including the
branching rule, are described in
\cite[Theorem~5.8, Section~6, and Corollary~7.1]{vershik-okounkov}.

\begin{lemma}\label{smooth:lem-dimension}
For each fixed integer \(d\ge1\) and all sufficiently large \(N\),
every irreducible constituent of the complexification of
\(\mathcal P_d^\perp\) has dimension at least
\[
 c_dN^{d/2},\qquad c_d=\frac1{2^dd!}.
\]
\end{lemma}

\begin{proof}
Fix a set of \(d\) vertices, and let \(G\cong S_{N-d}\) permute the
other vertices.  A \(G\)-orbit of matchings is determined by the
internal matching \(F\) on the fixed vertices: a permutation can
first align the outside partners of all unmatched fixed vertices
and then align the remaining outside edges.  The orbit indicator is
\[
 \prod_{e\in F}x_e
 \prod_{\substack{e\text{ on the fixed vertices}\\e\notin F}}(1-x_e).
\]
On expansion, a nonzero monomial uses only disjoint edges among
those \(d\) vertices, and therefore at most \(\lfloor d/2\rfloor\)
edge indicators.  All \(G\)-invariant functions consequently belong
to \(\mathcal P_d\).

If a Young diagram has first row of length at least \(N-d\), the
branching rule gives a \(G\)-fixed vector in its irreducible
representation: remove corner boxes until only a row of size
\(N-d\) remains.  Such a representation cannot occur in
\(\mathcal P_d^\perp\), since any such copy would provide a nonzero
\(G\)-fixed vector there.

A first column of length at least \(N-d\) similarly gives a vector
transforming by the sign character of \(G\).  No nonzero matching
function has this transformation law when \(N>2d\).  In fact,
every matching then contains an edge entirely outside the fixed
vertices; transposing its endpoints belongs to \(G\), fixes that
matching, and reverses the sign of the function.  The function must
vanish at every matching.  Thus a diagram occurring in
\(\mathcal P_d^\perp\) has both first row and first column shorter than
\(N-d\).

Let \(\lambda\) be any remaining diagram.  The product of its first
row and first column lengths is at least \(N\).  Transpose if
necessary so that its first row has length \(b\ge\sqrt N\);
transposition preserves the number of standard tableaux.  Retain
that row and a Young subdiagram of exactly \(d\) boxes below it.
Such a subdiagram exists because \(N-b\ge d\).

For large \(N\), we have \(b\ge2d\).  Put \(1,\ldots,d\) into the
first \(d\) positions of the top row.  Choose any \(d\) of the
remaining \(b\) labels \(d+1,\ldots,b+d\) for the lower boxes, fill
those boxes in a fixed standard order, and put the unused labels
increasingly into the rest of the top row.  Each lower box is in
one of the first \(d\) columns, so the resulting tableau is
standard.  The \(\binom bd\) choices give distinct tableaux of this
subdiagram.  Each extends to \(\lambda\) by adding the remaining
boxes along a fixed corner-addition chain with successively larger
labels.  The dimension \(f^\lambda\) therefore satisfies
\[
 f^\lambda\ge\binom bd
 \ge\frac{b^d}{2^dd!}
 \ge c_dN^{d/2},
\]
as required.
\end{proof}

\subsection{From moments to smoothing, and from products to polynomials}

\begin{proof}[Proof of Theorem~\ref{thm:matching-smoothing}]
Complexification does not change operator norms.  The operator
\(\mathcal A=T_t^*T_t\) commutes with vertex permutations, since
\(T_t\) intertwines their actions on matchings and partitions.
The projection \(Q=I-\Pi_d\) also commutes with permutations,
because \(\mathcal P_d\) is invariant.  The compression
\(B=Q\mathcal A Q\), acting on \(\mathcal P_d^\perp\), is therefore
positive semidefinite and equivariant.  Moreover,
\[
 \tr(B^2)\le\tr(\mathcal A^2),\qquad
 \normop{B}=\normop{T_tQ}^2.
\]
The first inequality is contraction of the Hilbert--Schmidt norm
under orthogonal compression and does not require \(Q\) and
\(\mathcal A\) to commute.

If \(B\ne0\), its largest-eigenvalue eigenspace is invariant and
contains an irreducible constituent.  Lemma~\ref{smooth:lem-dimension}
and the fourth-moment bound give
\[
 c_dN^{d/2}\normop{T_tQ}^4
 \le e^{2a_\ell}(N+1)^{b_\ell}.
\]
The same inequality holds if \(B=0\).  Choose \(d\) so that
\(d/2>b_\ell+4K\).  For all sufficiently large \(N\), this proves
\eqref{smooth:eq-smoothing}, uniformly in the sizes and the table.

To prove the polynomial assertion, fix a partition and a generating
product of edge indicators.  Repeated indicators may be removed,
and a product containing two distinct incident edges vanishes.
Otherwise the specified edges form a partial matching.  Let
\(m_{uv}\) be its type counts and let
\[
 a_u=2m_{uu}+\sum_{v\ne u}m_{uv}
\]
be the number of its endpoints in cell \(u\).  The probability
that a matching with table \(t\) contains these specified edges is
\begin{equation}\label{smooth:eq-containment}
 \frac{2^{\sum_um_{uu}}\prod_{u\le v}(t_{uv})_{m_{uv}}}
      {\prod_u(s_u)_{a_u}},
 \qquad
 (z)_j=z(z-1)\cdots(z-j+1),\quad (z)_0=1.
\end{equation}
Indeed, apply a uniform cell-preserving vertex permutation to any
one matching with table \(t\).  This is uniform over all such
matchings by transitivity.  The denominator counts injections of
the specified endpoint lists, while the numerator chooses distinct
edges of each type and the two orientations of each internal edge.
The denominator is positive because \(a_u\le s_u\).
If \(t_{uv}<m_{uv}\) for any type, its falling factorial is zero,
as required.

The numerator of \eqref{smooth:eq-containment} has total degree equal to
the number of specified edges.  Equation~\eqref{smooth:eq-table-degrees}
expresses every \(t_{uu}\) as an affine function of the intercell
counts.  Thus the containment probability is a polynomial of the
required degree in those counts.  By linearity this holds for
\(\Pi_df\); its expansion in edge products is fixed independently
of \(t\), so the resulting polynomial works simultaneously for all
admissible tables.
\end{proof}


With three cells, the table is determined by its three intercell edge counts.

\begin{corollary}[Three-cell smoothing]\label{cor:three-cell}
For every fixed $K>0$, there are an integer $d\ge1$ and $N_0$ with the following property. Let $N\ge N_0$ be even, and let $a,b,c$ be positive integers summing to $N$. For a uniform ordered partition into cells of these sizes, prescribe $x,y,z$ edges between the first and second, first and third, and second and third cells, respectively. Assume that $x,y,z$ are nonnegative integers and that
\[
 a-x-y,\qquad b-x-z,\qquad c-y-z
\]
are nonnegative even integers. Let $T_{x,y,z}$ average functions on perfect matchings subject to these counts, and let $\Pi_d$ project onto the edge-product span of degree at most $d$, using uniform measures as above. Then
\[
 \normop{T_{x,y,z}(I-\Pi_d)}\le N^{-K}
\]
uniformly in all the sizes and admissible triples. For every fixed partition and matching function $f$, the values $T_{x,y,z}\Pi_df$ are given by one polynomial of total degree at most $d$ in $(x,y,z)$ across all admissible triples.
\end{corollary}

\begin{proof}
Apply Theorem~\ref{thm:matching-smoothing} with $\ell=3$, intercell counts $t_{12}=x$, $t_{13}=y$, $t_{23}=z$, and internal counts $(a-x-y)/2$, $(b-x-z)/2$, $(c-y-z)/2$.
\end{proof}
